\section{Introduction}

\subsection{From Fourier optics to angle-resolved operator metasurfaces}

Fourier optics provides the formal language for interpreting imaging systems and linear spatial operators in the frequency domain, where filtering and differentiation are encoded through transfer functions rather than through pointwise geometric manipulations \cite{goodman2017}. In metasurface-enabled analog photonics, this viewpoint has become especially powerful because a subwavelength structured interface can be designed to emulate a desired optical operator within an extremely compact footprint. Recent demonstrations of mathematical metasurfaces and nonlocal operator devices have shown that differentiation, edge enhancement, and polarization-dependent image processing can be implemented by engineering the angular response of nanostructured optical layers \cite{silva2014,cordaro2019,kwon2020,cotrufo2023optica}.

For angle-resolved Fourier-operator metasurfaces, however, the relevant design object is not merely a scalar transmission efficiency or a single transfer-function slice at normal incidence. The device must realize a structured optical response that varies jointly with incidence angle, wavelength, and polarization. This coupled dependence is not a secondary implementation detail; it is the mechanism by which operator fidelity, angular bandwidth, and polarization selectivity are actually realized. Accordingly, a satisfactory inverse-design framework should be able to represent and manipulate the response field itself.

\subsection{Why inverse design remains difficult}

Despite rapid progress in metasurface inverse design, freeform Fourier-operator devices remain challenging because the structure-to-response map is high-dimensional, strongly nonlinear, and frequently non-unique. A large number of distinct binary patterns can yield similar coarse operator behavior while differing substantially in spectral robustness, angular stability, or polarization leakage. Brute-force scanning is therefore prohibitively inefficient, and purely deterministic regression tends to average over valid design modes instead of preserving solution diversity \cite{molesky2018,liu2018,so2019}.

The challenge becomes more pronounced when the target is defined over a multi-axis response field. Conventional design workflows often reduce the problem to one operating wavelength, one incidence condition, or one target task at a time. Such simplifications are useful for isolated demonstrations, but they obscure the central question addressed in this work: whether angle-resolved Fourier-operator metasurface design can be cast as a unified response-space inverse problem whose different practical instantiations emerge from one common formulation.

\subsection{Central thesis and scope of this paper}

The central thesis of this manuscript is that physics-consistent generative inverse design offers a natural framework for this unification. We treat the desired metasurface behavior as a target optical response field $\Rtar(\lam,\ang,p)$ defined over wavelength, incidence angle, and polarization channel. The inverse problem is then to generate one or more physically plausible structures whose verified response matches this target field under full-wave or \rcwa-based evaluation. Under this formulation, second-order differentiation, polarization-multiplexed image processing, and spectrally selective operator behavior are not separate ``tasks'' but different constraints on the same unified response object.

This narrative deliberately avoids framing the work as a platform paper built around a growing library of independent demonstrations. Instead, the paper is organized around a single design paradigm: angle--frequency--polarization response-space engineering enforced through a physics-consistent generative loop. The sections that follow formalize this viewpoint, describe the generative framework, and use representative results to argue that unification is the main scientific contribution.
